\subsection{The current sparse tools do not reproduce a positive target-score change}
\label{sec:sparse-effect}

Restoring the complete hidden-state difference increases target-answer scores, but the increase does not show that a fitted sparse reconstruction produces the same score change. We compare complete and sparse interventions at one matched internal object: the feed-forward output at layer 17 for Gemma, layer 14 for Qwen, and layer 16 for LLaVA. For each model, the experiment measures the feed-forward output in the complete-image and answer-region-corrupted runs, reconstructs the output difference with the sparse tool, and writes each candidate difference back at the same position. The combined intervention uses 32 positions sampled uniformly across the visual tokens and the four text positions immediately preceding the answer. Keeping the examples, answer readout, and write locations fixed separates the complete change, sparse reconstruction, and reconstruction residual. We first qualify reconstruction accuracy and then evaluate how inserting or deleting each component changes the complete-answer log probability.

\subsubsection{The sparse tools reconstruct the relevant change poorly}

The fitted sparse tools do not accurately reconstruct the change caused by evidence corruption in the compact-evidence questions studied here. At the primary $K=128$ budget, Equations~\ref{eq:relative-reconstruction-error} and~\ref{eq:reconstruction-direction-similarity} give mean answer-position errors of $1.756$ (Gemma, $n=60$), $1.872$ (Qwen, $n=60$), and $1.617$ (LLaVA, $n=59$), with corresponding direction similarities of $.227$, $.221$, and $.199$. One LLaVA image is excluded from both image-level means because its complete change is zero at all four positions. We therefore interpret the following results as measurements of what each fitted tool reveals, not as direct measurements of a complete sparse mechanism in the base model.

Even under that restricted interpretation, the comparison itself is well controlled. Each model uses the same 240 confirmation images, the same complete-answer log-probability readout, and the same write location for the complete change, sparse reconstruction, and reconstruction residual. At most 128 nonnegative feature activations per token form the main sparse reconstruction; fixed budgets of 16, 32, 64, and 128 are all retained. Original-state writeback, vector decomposition, joint writeback, random-direction norm matching, repeated forward passes, and agreement with the native-channel answer baselines all pass their numerical checks.

\subsubsection{Sparse and remaining components change target-answer scores differently}

Table~\ref{tab:same-object-components} reports the combined visual- and answer-position interventions at the main 128-feature budget. Positive restoration means that inserting a component from the complete-image run increases the correct answer's complete-sequence log probability in the corrupted run. Positive deletion means that removing the corresponding complete--corrupted difference decreases the same sequence score.

\begin{samepage}
\begin{center}
\begin{minipage}{\linewidth}
\captionsetup{type=table,hypcap=false}
\caption{\textbf{The selected sparse reconstruction does not produce a reliable positive target-score change.} Results combine 32 uniformly sampled visual positions with the four text positions immediately preceding the answer and use at most 128 active features per token on the same 240 images per model. Each cell gives the mean change in complete-answer log probability and an uncorrected 95\% cluster-bootstrap interval. Sparse is the selected reconstruction; residual is the reconstruction residual. $\ast$ marks results that remain different from zero after joint correction of 72 tests.}
\label{tab:same-object-components}
\centering
\small
\setlength{\tabcolsep}{2.2pt}
\begin{tabular*}{\linewidth}{@{\extracolsep{\fill}}lcccc@{}}
\toprule
Model & Sparse restore & Residual restore & Sparse delete & Residual delete \\
\midrule
Gemma & \shortstack{$-0.428^{\ast}$\\$[-0.599,-0.260]$} & \shortstack{$0.516^{\ast}$\\$[0.281,0.760]$} & \shortstack{$-0.540^{\ast}$\\$[-0.706,-0.372]$} & \shortstack{$0.314$\\$[0.110,0.515]$} \\
Qwen & \shortstack{$-0.005$\\$[-0.028,0.017]$} & \shortstack{$0.022$\\$[0.002,0.043]$} & \shortstack{$-0.008$\\$[-0.026,0.010]$} & \shortstack{$-0.007$\\$[-0.028,0.012]$} \\
LLaVA & \shortstack{$-0.017$\\$[-0.033,-0.002]$} & \shortstack{$0.000$\\$[-0.017,0.018]$} & \shortstack{$0.003$\\$[-0.008,0.015]$} & \shortstack{$0.010$\\$[-0.006,0.025]$} \\
\bottomrule
\end{tabular*}
\end{minipage}
\end{center}
\end{samepage}


No positive sparse-component score change survives the joint correction in any model. Gemma instead shows negative sparse restoration and deletion changes, while restoring the reconstruction residual increases the sequence score. Qwen and LLaVA show no component-level change after the same correction. The specificity controls reinforce this boundary: no sparse reconstruction outperforms equal-norm random directions or a wrong insertion position, and Gemma's combined-position sparse restoration is significantly weaker than both controls.

The four fixed feature budgets do not reveal a hidden positive pattern. Their mean sparse-restoration changes reverse direction across budgets and remain small for Qwen and LLaVA; Gemma's answer-position change remains negative from 16 through 128 features. These curves are descriptive because the primary joint tests use the frozen 128-feature setting. The conclusion is therefore about the current tools: the tested sparse reconstructions do not produce a reliable positive target-score change under restoration or deletion. This failure does not imply that the base models lack sparse computation, because reconstruction error may prevent the fitted interfaces from reproducing the complete internal change. Before removing the fitted tools from the next causal test, we examine what their sparse descriptions display across models and why those descriptions cannot by themselves establish a path.

\FloatBarrier
